\documentclass[12pt,a4paper]{article}
\usepackage[utf8]{inputenc}
\usepackage[T1]{fontenc}
\usepackage[french]{babel}
\usepackage{amsmath, amssymb}
\usepackage{graphicx}
\usepackage{geometry}
\usepackage{fancyhdr}
\usepackage{listings}
\usepackage{xcolor}
\usepackage{caption}
\usepackage{hyperref}
\usepackage{booktabs}
\usepackage{float}
\usepackage{enumitem}

\geometry{margin=2.5cm}

% Style MATLAB pour les listings
\definecolor{matlabgreen}{rgb}{0.13,0.55,0.13}
\definecolor{matlabblue}{rgb}{0.0,0.0,0.8}
\definecolor{matlabgray}{rgb}{0.5,0.5,0.5}
\lstdefinestyle{matlab}{
    language=Matlab,
    basicstyle=\ttfamily\small,
    keywordstyle=\color{matlabblue}\bfseries,
    commentstyle=\color{matlabgreen}\itshape,
    stringstyle=\color{red},
    numbers=left,
    numberstyle=\tiny\color{matlabgray},
    numbersep=8pt,
    frame=single,
    breaklines=true,
    tabsize=4,
    showstringspaces=false
}
\lstset{style=matlab}

\pagestyle{fancy}
\fancyhf{}
\lhead{APP Signal --- Exercice 6.5}
\rhead{Groupe G2D}
\cfoot{\thepage}

\title{
    \vspace{-1cm}
    \textbf{APP Signal --- Exercice 6.5} \\[0.3cm]
    \Large Sous-\'echantillonnage et Filtre Anti-Repliement \\[0.5cm]
    \large ISEP --- 2025--2026
}
\author{Groupe G2D}
\date{Avril 2026}

\begin{document}

\maketitle
\thispagestyle{fancy}
\tableofcontents
\newpage

% ===========================================================================
\section{Introduction}
% ===========================================================================

Cet exercice porte sur le \textbf{sous-\'echantillonnage} d'un signal audio de
$F_e = 44\,100$~Hz vers $F_e' = 4\,000$~Hz, avec un rapport signal \`a bruit
sup\'erieur \`a 40~dB. Il met en \'evidence le ph\'enom\`ene de
\textbf{repliement spectral} (aliasing) et l'importance du \textbf{filtre
anti-repliement} avant toute d\'ecimation.

L'ensemble du code est d\'evelopp\'e en MATLAB, \textbf{sans Signal Processing
Toolbox}, en utilisant uniquement des fonctions de base (FFT, \texttt{filter},
\texttt{interp1}).

% ===========================================================================
\section{Th\'eorie}
% ===========================================================================

\subsection{Th\'eor\`eme de Shannon}

Un signal \'echantillonn\'e \`a $F_e$ ne peut repr\'esenter fid\`element que
les composantes fr\'equentielles inf\'erieures \`a $F_e/2$ :
\begin{equation}
    F_e \geq 2\,f_{\max}
\end{equation}

Lors d'une d\'ecimation vers $F_e' < F_e$, toute composante au-dessus de
$F_e'/2$ se \textbf{replie} dans la bande $[0,\; F_e'/2]$, cr\'eant des
distorsions irr\'eversibles.

\subsection{Filtre anti-repliement}

Pour \'eviter le repliement, on applique un filtre passe-bas de fr\'equence de
coupure $f_c = F_e'/2$ \textbf{avant} la d\'ecimation. L'att\'enuation en bande
coup\'ee doit \^etre suffisante pour garantir le SNR souhait\'e.

% ===========================================================================
\section{Question 1 --- Chargement et affichage du signal}
% ===========================================================================

Le fichier \texttt{GardeMontante.wav} est un signal audio st\'er\'eo
\'echantillonn\'e \`a 44\,100~Hz. On le convertit en mono :

\begin{lstlisting}
[x, Fe] = audioread('GardeMontante.wav');
x = mean(x, 2);  % stereo -> mono
\end{lstlisting}

La densit\'e spectrale de puissance (DSP) est estim\'ee par FFT via la fonction
\texttt{myDSPEstimation} :

\begin{equation}
    \text{DSP}(f) = 10\log_{10}\!\left(\left|\frac{X(f)}{N}\right|^2\right) \quad \text{(dB)}
\end{equation}

\textbf{Observation} : le spectre s'\'etend au-del\`a de 2\,000~Hz ($= F_e'/2$),
ce qui annonce un repliement si l'on d\'ecime directement.

% ===========================================================================
\section{Question 2 --- D\'ecimation sans filtre}
% ===========================================================================

On r\'e-\'echantillonne directement le signal \`a $F_e' = 4\,000$~Hz par
interpolation lin\'eaire :

\begin{lstlisting}
Fe2 = 4000;
Te2 = 1 / Fe2;
t2  = 0 : Te2 : t(end);
xl2 = interp1(t, x, t2, 'linear');
\end{lstlisting}

Le facteur de d\'ecimation est :
\begin{equation}
    D = \frac{F_e}{F_e'} = \frac{44\,100}{4\,000} \approx 11
\end{equation}

\textbf{Observation} : la DSP du signal d\'ecim\'e montre un \textbf{repliement
spectral} clair. \`A l'\'ecoute, des artefacts sont audibles (distorsions,
bruits parasites).

% ===========================================================================
\section{Question 3 --- Conception du filtre anti-repliement}
% ===========================================================================

\subsection{Gabarit du filtre}

\begin{table}[H]
    \centering
    \begin{tabular}{lll}
        \toprule
        \textbf{Param\`etre} & \textbf{Valeur} & \textbf{Justification} \\
        \midrule
        Type     & Passe-bas FIR  & Phase lin\'eaire \\
        $F_e$    & 44\,100~Hz     & Fr\'equence du signal original \\
        $F_{\text{pass}}$ & 1\,800~Hz & Bande utile \`a conserver \\
        $F_{\text{stop}}$ & 2\,000~Hz & $= F_e'/2$, au-del\`a $\rightarrow$ repliement \\
        $A_{\text{stop}}$ & 40~dB     & SNR $> 40$~dB impos\'e \\
        $A_{\text{pass}}$ & 0,1~dB    & Ondulation minimale \\
        \bottomrule
    \end{tabular}
    \caption{Gabarit du filtre anti-repliement.}
\end{table}

\subsection{M\'ethode de Kaiser}

Le filtre est con\c{c}u par \textbf{fen\^etrage de Kaiser} (impl\'ementation
manuelle, sans toolbox). Le param\`etre $\beta$ est calcul\'e en fonction de
l'att\'enuation souhait\'ee :

\begin{equation}
    \beta = \begin{cases}
        0{,}1102\,(A - 8{,}7) & \text{si } A > 50 \\
        0{,}5842\,(A-21)^{0,4} + 0{,}07886\,(A-21) & \text{si } 21 \leq A \leq 50 \\
        0 & \text{sinon}
    \end{cases}
\end{equation}

L'ordre du filtre est estim\'e par :
\begin{equation}
    N = \left\lceil \frac{A - 7{,}95}{2{,}285 \cdot 2\pi \cdot \Delta f} \right\rceil
    \quad \text{avec } \Delta f = \frac{F_{\text{stop}} - F_{\text{pass}}}{F_e}
\end{equation}

La r\'eponse impulsionnelle est :
\begin{equation}
    h(n) = w_{\text{Kaiser}}(n) \cdot 2f_c \cdot \text{sinc}\!\left(2f_c\!\left(n - \frac{N}{2}\right)\right)
\end{equation}

avec la fen\^etre de Kaiser :
\begin{equation}
    w(n) = \frac{I_0\!\left(\beta\sqrt{1 - \left(\frac{2n}{N} - 1\right)^2}\right)}{I_0(\beta)}
\end{equation}

o\`u $I_0$ est la fonction de Bessel modifi\'ee de premi\`ere esp\`ece d'ordre z\'ero.

\begin{lstlisting}
[h, Nfilt] = myKaiserFilter(Fe, 1800, 2000, 40);
% Nfilt ~ 80-100 (ordre estime automatiquement)
\end{lstlisting}

\textbf{V\'erification} : la r\'eponse en fr\'equence montre une att\'enuation
$> 40$~dB au-del\`a de 2\,000~Hz et un gain $\approx 0$~dB en bande passante.
La phase est lin\'eaire (filtre FIR \`a coefficients sym\'etriques).

% ===========================================================================
\section{Question 4 --- Filtrage du signal}
% ===========================================================================

Le filtre est appliqu\'e au signal original :

\begin{lstlisting}
xl3 = filter(h, 1, x);
\end{lstlisting}

\textbf{Observation} : la DSP du signal filtr\'e montre que les composantes
au-dessus de 2\,000~Hz sont att\'enu\'ees de plus de 40~dB. \`A l'\'ecoute,
le signal est l\'eg\`erement plus sourd (perte des hautes fr\'equences) mais
sans artefacts.

% ===========================================================================
\section{Question 5 --- D\'ecimation du signal filtr\'e}
% ===========================================================================

On d\'ecime le signal filtr\'e :

\begin{lstlisting}
xl4 = interp1(t, xl3, t2, 'linear');
\end{lstlisting}

\textbf{Observation} : la DSP du signal d\'ecim\'e ne pr\'esente \textbf{plus de
repliement}. Le spectre est propre dans la bande $[0,\; 2\,000]$~Hz. \`A
l'\'ecoute, la qualit\'e est bonne malgr\'e la bande passante r\'eduite.

% ===========================================================================
\section{Comparaison des r\'esultats}
% ===========================================================================

\begin{table}[H]
    \centering
    \begin{tabular}{lccc}
        \toprule
        \textbf{Signal} & \textbf{$F_e$ (Hz)} & \textbf{Qualit\'e} & \textbf{Repliement} \\
        \midrule
        Original                 & 44\,100 & Excellente              & Non \\
        D\'ecim\'e sans filtre   & 4\,000  & D\'egrad\'ee            & \textbf{Oui} \\
        D\'ecim\'e avec filtre   & 4\,000  & Bonne (bande limit\'ee) & \textbf{Non} \\
        \bottomrule
    \end{tabular}
    \caption{Comparaison des trois signaux.}
\end{table}

% ===========================================================================
\section{Architecture logicielle}
% ===========================================================================

\begin{table}[H]
    \centering
    \begin{tabular}{ll}
        \toprule
        \textbf{Fichier} & \textbf{R\^ole} \\
        \midrule
        \texttt{main.m}             & Point d'entr\'ee, ajoute les paths \\
        \texttt{exercice\_6\_5.m}   & Script complet (Q1 \`a Q5, affichages) \\
        \texttt{myDSPEstimation.m}  & Estimation de la DSP par FFT \\
        \texttt{myKaiserFilter.m}   & Filtre FIR passe-bas (fen\^etre de Kaiser) \\
        \texttt{myDecimation.m}     & D\'ecimation avec/sans filtre anti-repliement \\
        \texttt{test\_*.m}          & Tests unitaires (10 cas par fonction) \\
        \bottomrule
    \end{tabular}
    \caption{Organisation des fichiers MATLAB.}
\end{table}

% ===========================================================================
\section{Conclusion}
% ===========================================================================

Cet exercice illustre concr\`etement le \textbf{th\'eor\`eme de Shannon} et ses
cons\'equences pratiques. La d\'ecimation directe (sans filtrage pr\'ealable)
provoque un repliement spectral audible, tandis que l'application d'un filtre
FIR anti-repliement con\c{c}u par fen\^etrage de Kaiser \'elimine ce
ph\'enom\`ene tout en pr\'eservant un SNR $> 40$~dB.

Points cl\'es mis en pratique :
\begin{itemize}
    \item Estimation de la DSP par FFT
    \item Conception d'un filtre FIR sans toolbox (formule de Kaiser + sinc)
    \item V\'erification par r\'eponse en fr\'equence et tests unitaires
    \item D\'ecimation par interpolation lin\'eaire (\texttt{interp1})
\end{itemize}

Le code est enti\`erement modulaire (3 fonctions r\'eutilisables + 30 tests
unitaires) et ne n\'ecessite aucune toolbox payante.

\end{document}
